\documentclass[../../../include/open-logic-section]{subfiles}

\begin{document}

\olfileid{sth}{spine}{recursion}
\olsection{د نامتناهيت هاخوا بازګښت قضيې}

خو لومړے کار دا دے چې تاييد کړو
\olref[valpha]{defValphas} يو بريالے تعريف دے. په عمومي توګه، ثابتول
پکار دي چې د نامتناهيت هاخوا په بازګښت د تعريف د وړاندې کولو هره هڅه
بريالۍ کېږي. د دې برخې مقصد همدا دے.

\emph{خبردار: دا مواد پېچلي دي}. خو عمومي مطلب ئې ساده دے:
د نامتناهيت هاخوا استقرا او تعويض په ګډه د نامتناهيت هاخوا بازګښت د
څو بڼو جواز تضمينوي.\footnote{يوه يادونه: ټول فارمولونه او ترمونه
پارامترونه لرلے شي، مګر که په ښکاره بل ډول ويل شوي وي.}

\begin{defn}
	فرض کړئ $\tau(x)$ ترم دے، $f$ تابعه ده، او $\alpha$ ترتيبي عدد دے.
وايو چې $f$ يو \emph{$\alpha$-تقريب} د $\tau$ دپاره دے، که او يوازې که
دواړه $\dom{f} = \alpha$ او $(\forall \beta \in \alpha)f(\beta) = \tau(\funrestrictionto{f}{\beta})$ وي.
\end{defn}

\begin{lem}[محدود بازګښت]\ollabel{transrecursionfun}
د هر ترم $\tau(x)$ او هر ترتيبي عدد $\alpha$ دپاره، يو يوازينے
$\alpha$-تقريب د $\tau$ دپاره شته.
\end{lem}
\begin{proof}
و به ښيو چې د هر $\gamma \leq \alpha$ دپاره يو يوازينے
$\gamma$-تقريب شته.

لومړے يکتايي ثابتوو. فرض کړئ $g$ او $h$ په ترتيب سره $\gamma$- او
$\delta$-تقريبونه دي. د دواړو تابعو په ورننوتونکو ترتيبي عددونو باندې د نامتناهيت هاخوا استقرا ښيي چې
$g(\beta) = h(\beta)$ د هر $\beta \in \dom{g} \cap \dom{h} =
\gamma \cap \delta = \min(\gamma, \delta)$ دپاره دي. نو زموږ تقريبونه،
که شتون ولري، يوازيني دي او په ټولو ګډو قيمتونو سره سمون لري.

د شتون د ثابتولو دپاره اوس د نامتناهيت هاخوا ساده استقرا
(\olref[ordinals][opps]{simpletransrecursion}) په هغو ترتيبي عددونو
کاروو چې $\delta \leq \alpha$ وي.

تشه تابعه په ښکاره $\emptyset$-تقريب ده.

که $g$ يو $\gamma$-تقريب وي، نو $g \cup
\{\tuple{\gamma, \tau(g)}\}$ يو $\ordsucc{\gamma}$-تقريب دے.

که $\gamma$ حدي ترتيبي عدد وي او $g_\delta$ يو
$\delta$-تقريب د هر $\delta < \gamma$ دپاره وي، نو $g = \bigcup_{\delta \in \gamma} g_\delta$ واخلو.
دا تابعه ده، ځکه چې بېل $g_\delta$ ګانې په ټولو ګډو قيمتونو سره سمون
لري. او که $\delta \in \gamma$ وي، نو $g(\delta) = g_{\ordsucc{\delta}}(\delta) =
\tau(\funrestrictionto{g_{\ordsucc{\delta}}}{\delta}) =
\tau(\funrestrictionto{g}{\delta})$ دے.

په دې سره د نامتناهيت هاخوا استقرا ثبوت بشپړېږي.
\end{proof}

که ځان ته د تابعې پر ځای د يوه \emph{ترم} د تعريفولو اجازه ورکړو، نو
د پخوانۍ نتيجې نه د $\alpha$ حد لرې کولے شو. د لاندې نتيجې په بيان او
ثبوت کښې، کله چې $\sigma$ ترم وي، نو
$\funrestrictionto{\sigma}{\alpha} = \Setabs{\tuple{\beta,
\sigma(\beta)}}{\beta \in \alpha}$ اخلو.

\begin{thm}[عمومي بازګښت]\ollabel{transrecursionschema}
د هر ترم $\tau(x)$ دپاره يو ترم
$\sigma(x)$ په ښکاره تعريفولے شو، داسې چې
$\sigma(\alpha) = \tau(\funrestrictionto{\sigma}{\alpha})$ د هر
ترتيبي عدد~$\alpha$ دپاره وي.
\end{thm}

\begin{proof}
د هر $\alpha$ دپاره د \olref{transrecursionfun} له مخې يو يوازينے
$\alpha$-تقريب $f_\alpha$ د~$\tau$ دپاره شته. $\sigma(\alpha)$ د
$f_{\ordsucc{\alpha}}(\alpha)$ په توګه تعريف کړو. اوس:
	\begin{align*}
		\sigma(\alpha) &=
		f_{\ordsucc{\alpha}}(\alpha) \\&=
		\tau(\funrestrictionto{f_{\ordsucc{\alpha}}}{\alpha}) \\&=
		\tau(\Setabs{\tuple{\beta, f_{\ordsucc{\alpha}}(\beta)}}{\beta \in \alpha}) \\&=
		\tau(\Setabs{\tuple{\beta, f_{\ordsucc{\beta}}(\beta)}}{\beta \in \alpha})\\&=
		\tau(\funrestrictionto{\sigma}{\alpha})
	\end{align*}
په ياد ولرئ چې $f_{\ordsucc{\beta}}(\beta) = f_{\ordsucc{\alpha}}(\beta)$
د ټولو $\beta < \alpha$ دپاره دي، لکه په \olref{transrecursionfun} کښې.
\end{proof}
\noindent
وګورئ چې \olref{transrecursionschema} يوه \emph{شېما} ده. دا مهمه ده چې
د $\sigma$ څخه د تابعې د تعريفولو تمه نۀ شو لرلے، يعنې د يوه ټاکلي
ډول \emph{سټ} د تعريفولو، ځکه چې بيا $\dom{\sigma}$ به د ټولو ترتيبي
عددونو سټ وي، چې د بورالي--فورتي له تناقض سره په ټکر کښې دے
(\olref[ordinals][basic]{buraliforti}).

خو لا ښودل پکار دي چې \olref{transrecursionschema}
زموږ د $V_\alpha$ ګانو تعريف توجيه کوي. ښايي دا سمدستي څرګند نۀ وي؛
خو د نامتناهيت هاخوا بازګښت په يوه وروستۍ ساده بڼه به څرګند شي.

\begin{thm}[ساده بازګښت]\ollabel{simplerecursionschema}
د هر دوو ترمونو $\tau(x)$ او $\theta(x)$ او هر سټ $A$ دپاره، يو ترم
$\sigma(x)$ په ښکاره تعريفولے شو چې:
\begin{align*}
	\sigma(\emptyset) &= A\\
	\sigma(\ordsucc{\alpha}) &= \tau(\sigma(\alpha)) &&
		\text{د هر ترتيبي عدد دپاره }\alpha\\
	\sigma(\alpha) &= \theta(\ran{\funrestrictionto{\sigma}{\alpha}})&&
	\text{کله چې }\alpha\text{ حدي ترتيبي عدد وي}
\end{align*}
\end{thm}

\begin{proof}
لومړے ترم $\xi(x)$ داسې تعريفوو:
%t $\xi(x) = A$ if $x$ is not a function whose domain is an ordinal;
%otherwise let:
\[
	\xi(x) =
	\begin{cases}
			A &\text{که $x$ تابعه نۀ وي، يا ئې د تعريف}\\
			&\hspace{1em}\text{ساحه ترتيبي عدد نۀ وي يا تشه وي؛ په بل صورت:}\\
			\tau(x(\alpha)) & \text{که $\dom{x} = \ordsucc{\alpha}$ وي}\\
			\theta(\ran{x}) & \text{که $\dom{x}$ حدي ترتيبي عدد وي}
	\end{cases}
\]
د \olref{transrecursionschema} له مخې يو ترم $\sigma(x)$ شته چې
$\sigma(\alpha) = \xi(\funrestrictionto{\sigma}{\alpha})$ د هر
ترتيبي عدد $\alpha$ دپاره وي؛ همدارنګه، $\funrestrictionto{\sigma}{\alpha}$
د $\alpha$ د تعريف ساحې لرونکې تابعه ده. ښيو چې $\sigma$ مطلوب خاصيتونه
لري، د نامتناهيت هاخوا ساده استقرا په وسيله
(\olref[ordinals][opps]{simpletransrecursion}). (د ژباړې يادونه
(OLSTH-005): د سرچينې په ښودل شوي تعريف کښې د تشې تابعې حالت نشته،
خو ورپسې ثبوت ئې بنسټيز قيمت کاروي. دلته په لومړي شرط کښې تشه
تعريف ساحه هم شامله شوې ده، څو هماغه بنسټيز قيمت ټاکل شوے وي.
نور تالي او حدي حالتونه او د ثبوت ټولې معادلې ساتل شوې دي.)

لومړے، $\sigma(\emptyset) = \xi(\emptyset) = A$ دے.

بيا، $\sigma(\ordsucc{\alpha}) = \xi(\funrestrictionto{\sigma}{\ordsucc{\alpha}}) = \tau(\funrestrictionto{\sigma}{\ordsucc{\alpha}}(\alpha)) = \tau(\sigma(\alpha))$ دے.

په پای کښې، $\sigma(\alpha) = \xi(\funrestrictionto{\sigma}{\alpha}) = \theta(\ran{\funrestrictionto{\sigma}{\alpha}})$، کله چې $\alpha$ حدي عدد وي.
\end{proof}
\noindent
اوس، د \olref[valpha]{defValphas} د توجيه کولو دپاره يوازې $A
= \emptyset$ او $\tau(x) = \Pow{x}$ او $\theta(x) = \bigcup x$ واخلو.
بالاخره په دې سره د $V_\alpha$ ګانو تعريف توجيه کېږي!{}

\end{document}
